\chapter{Literature Review}

\section{Evolution of Ground Control Stations and MAVLink}
The trajectory of Ground Control Station (GCS) development is intimately intertwined with the evolution of the Unmanned Aerial Vehicles (UAVs) they control. In the nascent stages of autonomous flight, GCS ecosystems were strictly proprietary, siloed operations. Corporations such as AeroVironment or General Atomics developed monolithic software-hardware stacks—such as the advanced stations utilized for the MQ-1 Predator—which required dedicated server racks, highly specialized cryptographic radio arrays, and thousands of hours of rigid operator training. These systems functioned on closed, undocumented communication protocols.

The democratization of UAV technology, spearheaded by the open-source community around the Pixhawk hardware standard and the ArduPilot/PX4 software stacks, necessitated a standardized communication protocol. This requirement birthed \textbf{MAVLink (Micro Air Vehicle Link)}, an incredibly lightweight, header-only message marshaling library. Introduced by Lorenz Meier in 2009, MAVLink fundamentally revolutionized autonomous robotics. By providing a serialized, binary protocol designed specifically for resource-constrained environments (operating over inherently lossy radio-frequency serial links), it allowed diverse flight controllers and diverse GCS software to communicate seamlessly.

MAVLink's architecture involves packaging data payloads into highly efficient, C-struct-like packets. MAVLink version 1 encapsulates messages in 8-26 byte packets, while version 2 introduces packet signing, larger payload capacities, and extension fields. The universal adoption of this protocol effectively decoupled the GCS from the UAV hardware, paving the way for numerous generic, cross-compatible ground station applications.

\section{Critical Analysis of Dominant Open-Source GCS Landscapes}
Despite the protocol standardization achieved by MAVLink, the design philosophy behind prominent GCS software has remained heavily rooted in legacy aerospace engineering paradigms. A comprehensive review of the current market reveals two dominant monolithic entities, alongside a nascent push toward web-centric alternatives.

\subsection{Mission Planner (ArduPilot Ecosystem)}
Developed by Michael Oborne, Mission Planner \cite{missionplanner} is undeniably the most feature-complete GCS for the ArduPilot ecosystem. It provides exhaustive capabilities: from initial rudimentary firmware flashing and PID (Proportional-Integral-Derivative) loop tuning, to incredibly complex 3D grid mission generation for photogrammetry. 

However, structurally, Mission Planner constitutes a classic example of an "expert tool." Built utilizing the historical Microsoft .NET Framework Windows Forms (WinForms), its user interface is notoriously complex. Flight operators are presented with a dizzying array of nested menus, configuration trees, and dense tabular data. In a high-stakes, time-sensitive Search and Rescue (SAR) deployment, this density translates directly into increased cognitive friction. Furthermore, its native reliance on Windows APIs creates severe deployment limitations; running Mission Planner on Linux or macOS necessitates cumbersome, resource-heavy emulation layers such as Mono, which frequently introduce UI artifacting, serial port instability, and rendering latency.

\subsection{QGroundControl (QGC)}
Recognizing the multi-platform limitations of Mission Planner, the Dronecode Project \cite{qgroundcontrol} developed QGroundControl (QGC). QGC was architected utilizing the Qt cross-platform application framework (C++). This allowed for true native compilation across Windows, macOS, Linux, iOS, and Android. QGC successfully modernized the user interface aesthetic, adopting a more touch-friendly, map-centric paradigm.

Nevertheless, QGC's fundamental reliance on the C++/Qt ecosystem presents a formidable barrier to rapid development and customization. The compilation pipeline is extensively heavy and deeply entangled with embedded systems concepts. Consequently, modern software engineers—predominantly trained in agile, reactive web stacks—find adapting, styling, or extending QGC for highly niche, localized use-cases functionally prohibitive without mastering Qt's complex signaling and slots architecture. Adding a simple, custom SVG-based payload control widget to QGC typically requires days of C++ interface modifications and complete binary recompilations.

\subsection{Web-Based and Distributed Architectures}
In recent commercial applications, there has been a significant migration toward web-centric (HTML/CSS/JS) architectures. Platforms like DroneDeploy or typical fleet management solutions route UAV telemetry through LTE/5G uplinks to centralized cloud servers \cite{nextjs}. The GCS becomes a mere thin-client browser window. 

While this architecture guarantees supreme cross-platform compatibility and rapid, component-based UI development (utilizing libraries like React or Vue), it introduces a fatal flaw for remote deployments: \textbf{Internet Dependency}. Search and rescue operations frequently occur in disaster zones, deep wilderness, or offshore environments where cellular infrastructure is either nonexistent or profoundly degraded. Cloud-tethered web architecture is inherently incapable of fulfilling the requirements of an autonomous, localized, off-grid GCS.

\section{Identifying the Architectural Gap}
A synthesis of the existing literature clearly delineates an unresolved paradigm in autonomous vehicle control: A rigid dichotomy exists between highly reliable, offline, C++/C\# desktop applications possessing poor DX (Developer Experience) and complex UIs, versus modern, highly customizable web applications that hopelessly rely on continuous Internet connectivity.

\textbf{The defined gap:} There is an acute, pressing requirement for a hybridized Ground Control Station. Such a GCS must capture the radical UI customization, fluid animations, and vast open-source component ecosystem of modern web frameworks (React/Next.js) \cite{nextjs}\cite{electronjs}, while executing entirely locally, offline, and maintaining low-level, high-bandwidth USB serial connections to 915MHz MAVLink telemetry radios. The Sky ResQ Dashboard is explicitly engineered to occupy this vacant niche.
